\section{A structural reduction to at most two doubled pairs}\label{sec:four-structure}

Throughout this section, cycles in a loopless multigraph distinguish parallel edges.
A pair of parallel edges forms a cycle of length two.
We prove that a minimum counterexample to $c(G)\le 2^n+n$ is simple or has at most two doubled pairs.
Together with Section~\ref{sec:four-finite}, this leaves twelve finite classes.

\begin{lemma}[Gluing across two doubled pairs]\label{lem:four-glue}
For $i=1,2$, let $H_i$ be a connected four-regular multigraph on $n_i$ vertices.
Choose two parallel edges with distinct ends $u_i,v_i$ in $H_i$.
Assume that deleting the chosen pair leaves a connected graph $N_i$.
Delete the two chosen pairs and add two parallel edges from $u_1$ to $u_2$
and two parallel edges from $v_1$ to $v_2$.
Call the resulting graph $G$.
If $c(H_i)\le 2^{n_i}+n_i$ for both $i$, then $c(G)\le 2^{n_1+n_2}+n_1+n_2$.
The chosen pair is permitted to be part of a larger parallel class.
\end{lemma}

\begin{proof}
Put $A_i=2^{n_i}$ and let $a_i=2P_{u_iv_i}(N_i)$.
Let $b_i=c(N_i)$.
The elementary path-to-cycle-space injection gives
$2\le a_i\le 2\cdot2^{\beta(N_i)}=A_i$,
since $\beta(N_i)=n_i-1$.
Classifying cycles by the chosen pair gives $c(H_i)=a_i+b_i+1$.
Classifying cycles by the two new pairs gives
\[
 c(G)=a_1a_2+b_1+b_2+2.
\]
Indeed, a cycle that crosses between the shores uses one edge from each new pair,
and a terminal path in each $N_i$.
The other cycles lie inside a shore or are one of the two new digons.
It follows that
\[
 c(G)\le a_1a_2-a_1-a_2+A_1+A_2+n_1+n_2.
\]
Finally,
\[
 A_1A_2-a_1a_2
 =(A_1-a_1)A_2+a_1(A_2-a_2)
 \ge (A_1-a_1)+(A_2-a_2).
\]
This proves the result.
\end{proof}

\begin{lemma}[A parity bound]\label{lem:four-parity}
Let $G$ be a connected four-regular multigraph on $n$ vertices whose only parallel classes
are $k$ doubled pairs forming a matching.
Delete all $2k$ edges in these pairs to obtain $F$, and put $t=\operatorname{comp}(F)$.
If $k-t\ge2$, then
\[
 c(G)\le \frac{27}{32}2^n-1+k<2^n+n.
\]
\end{lemma}

\begin{proof}
Contract each component of $F$ to a vertex and represent each doubled pair by one edge.
The resulting connected multigraph $J$ has $t$ vertices and $k$ edges; loops are permitted.
Let $Z$ count even edge sets of $G$ that use at most one edge from each doubled pair.
For each choice of pairs used, parity completion inside $F$ is possible exactly when
the corresponding edge set of $J$ is even.
There are then $2^{\beta(F)}$ completions and two choices for each selected pair.
Since $\beta(F)=n-2k+t$,
\[
 Z=2^{n-2k+t}\sum_{S\in\mathcal C(J)}2^{|S|}.
\]
The usual parity filter, obtained by expanding the product and summing vertex signs, gives
\[
 Z=2^{n-2k}\sum_{\sigma\in\{-1,1\}^{V(J)}}
       \prod_{uv\in E(J)}(1+2\sigma_u\sigma_v).
\]
Fix a spanning tree of $J$.
For an upper bound on the absolute value of the sum, each non-tree edge contributes at most three.
The sum of the absolute tree products is $2\cdot4^{t-1}$:
choose a root sign and then sum independently over agreement or disagreement on each tree edge.
Consequently,
\[
 \frac{Z}{2^n}\le 2^{-2k}\,2\cdot4^{t-1}3^{k-t+1}
       =\frac32\left(\frac34\right)^{k-t}.
\]
Every cycle except a doubled-pair digon is one of the nonempty even sets counted by $Z$.
Therefore $c(G)\le Z-1+k$, as required.
\end{proof}

\begin{theorem}\label{thm:four-two-pairs}
If a counterexample to $c(G)\le2^n+n$ exists, a counterexample with the fewest vertices
is either simple or has exactly one or two doubled pairs.
In the latter case, deleting all doubled pairs leaves a connected graph.
\end{theorem}

\begin{proof}
The complement and doubled-chain reductions in Section~\ref{sec:four-regular}
show that a minimum counterexample has a Hamilton decomposition, no multiplicities above two,
and a matching of doubled pairs.
Let $F,J,k,t$ have the meanings used in Lemma~\ref{lem:four-parity}.
We may assume $k\ge1$.

Each doubled pair contributes one edge to each Hamilton cycle:
using both would form a two-edge component of a Hamilton cycle on at least three vertices.
Thus every vertex of $J$ has even degree.
The graph $J$ is connected and Eulerian.
The parity bound gives $k-t\le1$.

Suppose $t>1$.
Every vertex of $J$ has positive even non-loop degree.
If all these degrees were at least four, then $k\ge2t$, contradicting $k-t\le1$.
Thus some vertex has non-loop degree two.
Its two incident non-loop edges form a cut with connected shores.
To see that the complementary shore is connected, any component there must have a positive even
number of cut edges, and only two cut edges are available.

The corresponding cut of $G$ consists of two doubled pairs.
Their four endpoints are distinct because the doubled pairs form a matching.
On each shore, close its two boundary vertices with a new doubled pair to form $H_i$.
Each $H_i$ is a connected loopless four-regular graph with fewer vertices than $G$.
The shore remains connected when the new pair is deleted.
Minimality gives $c(H_i)\le2^{n_i}+n_i$.
Lemma~\ref{lem:four-glue} then gives the required bound for $G$, a contradiction.
Hence $t=1$, and $k-t\le1$ gives $k\le2$.
\end{proof}

\paragraph{Status.}
This reduction and the analytic cutoffs are mathematical arguments.
The remaining finite cases require a complete enumeration or a further argument.
An enumeration is being prepared in separate source files; its outcome is not assumed here.
